\documentclass[12pt,letterpaper]{article}

%%
%% Required packages:
%%
\usepackage{etex}
\usepackage{amsmath}
\usepackage{amssymb}
\usepackage{amstext}
\usepackage{amsthm}
\usepackage{fancyhdr,fancybox}
\usepackage{mathrsfs} % need for \mathscr command
\usepackage{multicol}
\usepackage[normalem]{ulem}
%\usepackage{pictex,pictexplus}
% \usepackage{pictexwd}
\usepackage{rotating}
\usepackage{url}
\usepackage{xfrac}
\usepackage{booktabs}
\usepackage{color,soul}
\usepackage{comment}

\usepackage{hyperref}
\hypersetup{
    colorlinks=true,     % false: boxed links; true: colored links
    linkcolor=blue,      % color of internal links
    citecolor=blue,      % color of links to bibliography
    filecolor=blue,      % color of file links
    urlcolor=blue        % color of external links
}


\usepackage{tikz}
\usetikzlibrary{backgrounds,calc}

%%
%% Common formatting commands
%%
\setlength{\parindent}{0in}
\setlength{\textwidth}{7in}
\setlength{\evensidemargin}{-0.25in}
\setlength{\oddsidemargin}{-0.25in}
\setlength{\parskip}{.5\baselineskip}
\setlength{\topmargin}{-0.5in}
\setlength{\textheight}{9in}

%               Problem and Part
\newcounter{problemnumber}
\newcounter{partnumber}
\newcounter{subpartnumber}
\newcommand{\Problem}{%
  \refstepcounter{problemnumber}%
  \setcounter{partnumber}{0}%
  \setcounter{subpartnumber}{0}%
  \item[\makebox{\hfil\textbf{\theproblemnumber.}\hfil}]%
  }
\newcommand{\InProblem}[2][1.6]{%
  \refstepcounter{problemnumber}%
  \setcounter{partnumber}{0}%
  \setcounter{subpartnumber}{0}%
  \textbf{\theproblemnumber.}\ \ \parbox[t]{#1in}{#2}%
}
\newcommand{\InSmallProblem}[1]{\InProblem[1.05]{#1}}
\newcommand{\NoProblem}[1]{%
  \mbox{\hfil\textbf{#1}\hfil}%
  }
\newcommand{\UnProblem}{%
  \refstepcounter{problemnumber}%
  \setcounter{partnumber}{0}%
  \setcounter{subpartnumber}{0}%
  \mbox{\hfil\textbf{\theproblemnumber.}\hfil}%
  }
\newcommand{\InFlexProblem}[1]{%
  \refstepcounter{problemnumber}%
  \setcounter{partnumber}{0}%
  \setcounter{subpartnumber}{0}%
  \textbf{\theproblemnumber.}\ \ #1%
}
%%  Part:
\newcommand{\Part}{%
  \renewcommand{\thepartnumber}{(\alph{partnumber})}%
  \refstepcounter{partnumber}
  \setcounter{subpartnumber}{0}%
  \item[(\alph{partnumber})]%
}
\newcommand{\InPart}[2][1.75]{%
  \renewcommand{\thepartnumber}{(\alph{partnumber})}%
  \refstepcounter{partnumber}%
  \setcounter{subpartnumber}{0}%
  (\alph{partnumber})\ \ \parbox[t]{#1in}{#2}%
}
\newcommand{\UnPart}{%
  \refstepcounter{partnumber}%
  \renewcommand{\thepartnumber}{(\alph{partnumber})}%
  \setcounter{subpartnumber}{0}%
  (\alph{partnumber})%
}
\newcommand{\InLargePart}[1]{\InPart[3.05]{#1}}
\newcommand{\InFullPart}[1]{\InPart[6.2]{#1}}
\newcommand{\InSmallPart}[1]{\InPart[1.05]{#1}}
%% SubPart:
\newcommand{\SubPart}{%
  \refstepcounter{subpartnumber}%
  \item[(\roman{subpartnumber})]%
  }
\newcommand{\InSubPart}[2][2.25]{%
  \stepcounter{subpartnumber}%
  (\roman{subpartnumber})\ \ \parbox[t]{#1in}{#2}%
}
\newcommand{\InSmallSubPart}[1]{\InSubPart[1.5]{#1}}
\newcommand{\InTinySubPart}[1]{%
  \stepcounter{subpartnumber}%
  (\roman{subpartnumber})\ \ \makebox{#1}%
}
\newcommand{\InMySubPart}[2][2.25]{%
  \stepcounter{subpartnumber}%
  \parbox[t]{#1in}{\makebox[0.3in][l]{(\roman{subpartnumber})}\ \ {#2}}%
}
\newcommand{\TrueFalse}{\textbf{T}\ \ \textbf{F}\ \ }
\newcommand{\TFPart}[1]{% 
  \stepcounter{partnumber}%
  \setcounter{subpartnumber}{0}%
  \item[(\alph{partnumber})] \TrueFalse\parbox[t]{5.5in}{#1}}

\newcommand{\Points}[1]{(#1 points) \ }
\newcommand{\Point}[1]{(#1 point) \ }


%% For Answers:
\newcounter{answernumber}
\newcommand{\Answer}{%
  \refstepcounter{answernumber}%
  \setcounter{partnumber}{0}%
  \setcounter{subpartnumber}{0}%
  \item[\makebox{\hfil\textbf{\theanswernumber.}\hfil}]%
  }


% Setting up the headers for the first page
%\chead{} % will default to what is typed in the file.
\fancyhead[L]{\textsc{Math 22a}}
\fancyhead[R]{\textsc{Fall 2024}}
\fancyfoot[R]{
	\vspace*{\fill}\footnotesize\textcopyright{} \emph{Harvard Math 22a}	
}
\renewcommand{\headrulewidth}{0pt}

%\fancyhead[C]{}
%	\chead{}	
%	\lhead{}
%	\rhead{}
%	\cfoot{}


% Putting copyright on the footer of each page
%\fancypagestyle{secondpage}{
%	\fancyhead[C]{TESTing again} % change this to be blank
%		%\chead{}	
%	\fancyhead[L]{bears} % change this to be blank
%	\fancyhead[R]{dogs} % change this to be blank
%	\fancyfoot[R]{ %keeps the footer the same
%		\vspace*{\fill}\footnotesize\textcopyright{} \emph{Harvard Math 22a}	
%	}
%}
%\renewcommand{\headrulewidth}{0pt}
%\pagestyle{fancy}
\pagestyle{empty}


%\lhead{}
%\rhead{}
%\rfoot{\vspace*{\fill}\footnotesize\textcopyright{} \emph{Harvard Math 22a}}
%\renewcommand{\headrulewidth}{0pt}
%\pagestyle{fancy}




%%
%% Common shortcut commands:
%%
\newcommand{\ds}{\displaystyle}
\newcommand{\sgn}{\operatorname{sgn}}
\renewcommand{\Pr}{\operatorname{P}}
\newcommand{\CPr}[3][{}]{\Pr_{\text{#1}}\left(#2\,|\,#3\right)}

\newcommand{\C}{\mathbb{C}}
\newcommand{\R}{\mathbb{R}}
\newcommand{\Z}{\mathbb{Z}}
\newcommand{\N}{\mathbb{N}}
\newcommand{\Q}{\mathbb{Q}}
\newcommand{\D}{\mathbb{D}}

\newcommand{\rref}{\operatorname{rref}}
\newcommand{\im}{\operatorname{im}}
\newcommand{\rank}{\operatorname{rank}}
\newcommand{\diag}{\operatorname{diag}}
\newcommand{\Span}{\operatorname{span}}
%\renewcommand{\vec}[1]{\mathbf{#1}}
\newcommand{\charpoly}{\mathscr{P}}
\newcommand{\nicefrac}[2]{\sfrac{#1}{#2}} % using xfrac package
\newcommand{\topology}{\mathcal{T}}
\newcommand{\Inn}{\operatorname{Inn}}

\newcommand{\blankbox}[2]{\fbox{\rule{#1}{0in}\rule{0in}{#2}}}

\newenvironment{myindentpar}[1]%
 {\begin{list}{}%
         {\setlength{\leftmargin}{#1}
          \setlength{\rightmargin}{\leftmargin}}%
         \item[]%
 }
 {\end{list}}
\newtheorem{theorem}{Theorem}
\newtheorem*{NStheorem}{Theorem}
\newtheorem{cor}[theorem]{Corollary}
\newtheorem*{claim}{Claim}
\newtheorem{defn}{Definition}

\newenvironment{amatrix}[1]{%
  \left(\begin{array}{@{}*{#1}{c}|c@{}}
}{%
  \end{array}\right)
}

%--------grstep
% For denoting a Gauss' reduction step.
% Use as: \grstep{\rho_1+\rho_3} or \grstep[2\rho_5 \\ 3\rho_6]{\rho_1+\rho_3}
\newcommand{\grstep}[2][\relax]{%
   \ensuremath{\mathrel{
       {\mathop{\longrightarrow}\limits^{#2\mathstrut}_{
                                     \begin{subarray}{l} #1 \end{subarray}}}}}}
\newcommand{\swap}{\leftrightarrow}



\chead{\Large\textbf{Proof Methods, \fbox\S 4,  5,  6}}

\begin{document}
\thispagestyle{fancy}

\begin{enumerate}
  \Problem Give a \emph{direct proof}\ of the following statements.
  \begin{enumerate}
    \Part  Recall that, for natural numbers $n$ and $k$ with $k\le n$,
    the binomial coefficient is defined as $\binom{n}{k} =
    \frac{n!}{k!\,(n-k)!}$.  Show that, if $n\ge2$ is an integer, that
    $n^2 = 2\binom{n}{2} + \binom{n}{1}$.
    \vfill

    \Part Let $a$ and $b$ be integers. If $a|b$, then
    $a|\left(3b^3-4b^2+18b+2a\right)$.  \vfill

    \Part If $x$ is an integer and $x^2|x$, then $x\in
    \{-1,0,1\}$. 
    \vfill

  \end{enumerate}

  \Problem Give a \emph{contrapositive proof}\ of the following statements. 
  \begin{enumerate}
    \Part Suppose $x$ and $y$ are integers.  If $3\not|\,xy$,
    then $3\not|\,x$ and $3\not|\,y$.
    \vfill

    \Part (Hammack Problem 5.6) Suppose $x$ is a real number.  If $x^3
    -x > 0$, then $x>-1$.
    \vfill

    \Part (Hammack Problem 5.7) Suppose $a$ and $b$ are integers.  If
    both $ab$ and $a+b$ are even, then both $a$ and $b$ are even.
    \vfill
  \end{enumerate}

\end{enumerate}
\newpage

Prove the following using \emph{proof by contradiction}.
\begin{enumerate}
  \Problem $\sqrt{2}$ is irrational.
  \vfill

  \Problem There are infinitely many prime numbers.
  \vfill

 % \Problem (Hammack Problem 6.8) Suppose $a^2+b^2 = c^2$ for integers
%  $a$, $b$ and $c$.  Then at least one of $a$ and $b$ must be even.
 % \vfill

  \Problem (Basically Hammack Problem 6.16)\footnote{%
    This is also called the AM-GM Inequality, as it says that the
    Arithmetic Mean $\frac{x+y}{2}$ is always greater than or equal to
    the Geometeric Mean $\sqrt{xy}$.}   % 
  Suppose $x$ and $y$ are non-negative real numbers.  Then
  \begin{equation*}
    \frac{x+y}{2} \ge \sqrt{xy}
  \end{equation*}
  or, equivalently, $x+y\ge 2\sqrt{xy}$.  
  \vfill
\end{enumerate}

\end{document}

\newpage
\chead{\Large\textbf{Proof Methods -- Solutions}}
\lhead{}
\rhead{}
\thispagestyle{fancy}

\begin{enumerate}
  \Answer Give a \emph{direct proof}\ of the following statements.
  \begin{enumerate}
    \Part  Recall that, for natural numbers $n$ and $k$ with $k\le n$,
    the binomial coefficient is defined as $\binom{n}{k} =
    \frac{n!}{k!\,(n-k)!}$.  Show that, if $n\ge2$ is an integer, that
    $n^2 = 2\binom{n}{2} + \binom{n}{1}$.
  \begin{proof}
      	This will be a direct proof. So, take $2 \binom{n}{2} + \binom{n}{1}$ and substituting in for the binomial expressions, we get:
    	$$2 \binom{n}{2} +  \binom{n}{1} = 2 \frac{n!}{2!\,(n-2)!} + \frac{n!}{1!\,(n-1)!}  = 2 \frac{n(n-1)}{2} + \frac{n}{1} =$$ $$ n(n-1) + n = n^2 - n + n = n^2$$
  \end{proof}

    \Part Let $a$ and $b$ be integers. If $a|b$, then
    $a|\left(3b^3-4b^2+18b+2a\right)$.  
    
    \begin{proof}
        	If $a|b$, then $ b = ak$ for some integer $k$. So then we can substitute in for $b$ in the expression and get:
    	$$3b^3-4b^2+18b+2a = 3(ak)^3 - 4 (ak)^2 + 18(ak) + 2a = a(3a^2k^3 - 4ak^2 + 18k + 2)$$
    	Note that $3a^2k^3 - 4ak^2 + 18k + 2$ is an integer, so $a|b$. 
    \end{proof}

    \Part If $x$ is an integer and $x^2|x$, then $x\in
    \{-1,0,1\}$. 
  \begin{proof}
      	If $x^2|x$, we have that $x = x^2 k$ for some integer $k$. We have two cases. If $x$ is not zero, then we can cancel an $x$ and get that $xk = 1$, so $x|1$, so $x = \pm 1$. The other case is if $x=0$. So $x \in \{ -1,0,1 \}$.  
  \end{proof}

  \end{enumerate}

  \Answer Give a \emph{contrapositive proof}\ of the following statements. 
  \begin{enumerate}
    \Part Suppose $x$ and $y$ are integers.  If $3\not|\,xy$,
    then $3\not|\,x$ and $3\not|\,y$.
  \begin{proof}
      	The contrapositive statement is: if $3|x$ or $3|y$, then $3|xy$. So assume without loss of generality $3|x$, so $x = 3k$ for some integer $k$. So $xy = (3k)y = 3(ky)$. As $ky$ is an integer, $3| xy$. 
  \end{proof}

    \Part (Hammack Problem 5.6) Suppose $x$ is a real number.  If $x^3
    -x > 0$, then $x>-1$.
  \begin{proof}
      	The contrapositive statement is: if $x \leq -1$, then $x^3 - x \leq 0$. (Note here that the opposite of greater than is less than or equal to) If $x \leq -1$, then $x^3 \leq x$, so $x^3 - x \leq 0$. 
  \end{proof}

    \Part (Hammack Problem 5.7) Suppose $a$ and $b$ are integers.  If
    both $ab$ and $a+b$ are even, then both $a$ and $b$ are even.
  \begin{proof}
      	The contrapositive statement is: if $a$ or $b$ is odd, then $a+b$ or $ab$ is odd. There are two cases here. The first case is if only one of them are odd. So assume without loss of generality that $a$ is odd and $b$ is even, so $a = 2k+1$ and $b = 2j$ for integers $j,k$. Then, $a+b = 2k + 1 + 2j = 2(k+j) + 1$, which is odd. 
    	
    	The second case is if both are odd. So $a = 2k+1$ and $b = 2j+1$ for integers $j,k$. Then, $ab = (2k+1)(2j+1) = 4jk + 2j + 2k + 1 = 2(2jk + j  +k) + 1$ is odd. 
  \end{proof}
  \end{enumerate}

\end{enumerate}
\newpage

Prove the following using \emph{proof by contradiction}.
\begin{enumerate}
  \Answer $\sqrt{2}$ is irrational.
  \begin{proof}
    	Assume towards a contradiction that $\sqrt{2}$ is rational. Then $\sqrt{2} = \frac{p}{q}$. We can assume that $p$ and $q$ have no common factors by reducing. So, squaring both sides gives us $2 = \frac{p^2}{q^2}$, so $2q^2 = p^2$, which means $p^2$ is even. However, this also means $p$ is even. So, $p= 2k$ for some integer $k$, so $p^2 = 4k^2 = 2q^2$. This gives us that $2k^2 = q^2$, so $q^2$ is even, so $q$ is even. However, this is a contradiction, as $p$ and $q$ are both even but have no common factors. 
  \end{proof}

  \Answer There are infinitely many prime numbers.
  \begin{proof}
    	Suppose towards a contradiction that there are finitely many prime numbers $p_1,...,p_n$, with $p_n$ the largest prime number. Consider the number $N = p_1...p_n + 1$. This number is larger than our largest prime, so it can't be prime, so it must be divisible by some $p_i$, a prime. However, if $p_i$ divides $N$ and $p_i $ divides $p_1 ... p_n$, then $p_i$ divides $N - p_1 ... p_n = 1$, which is a contradiction.  
  \end{proof}

%  \Answer (Hammack Problem 6.8) Suppose $a^2+b^2 = c^2$ for integers
%  $a$, $b$ and $c$.  Then at least one of $a$ and $b$ must be even.
%  \begin{proof}
%    	Suppose towards a contradiction that $a$ and $b$ are both odd. Then $a = 2k+1$ and $b = 2j+1$, so $a^2 + b^2 = (2k+1)^2 + (2j+1)^2 = 4k^2 + 4k + 1 + 4j^2 + 4j + 1 = c^2$. So we have $2(2k^2 + 2k + 2j^2 + 2j + 1) = c^2$. However, this means that $2| c^2$, but $4 \not|\, c^2$, which is a contradiction. 
%  \end{proof}

  \Answer (Basically Hammack Problem 6.16)\footnote{%
    This is also called the AM-GM Inequality, as it says that the
    Arithmetic Mean $\frac{x+y}{2}$ is always greater than or equal to
    the Geometeric Mean $\sqrt{xy}$.}   % 
  Suppose $x$ and $y$ are non-negative real numbers.  Then
  \begin{equation*}
    \frac{x+y}{2} \ge \sqrt{xy}
  \end{equation*}
  or, equivalently, $x+y\ge 2\sqrt{xy}$.  
  \begin{proof}
    	Suppose towards a contradiction that $x+y < 2\sqrt{xy}$. As $x$ and $y$ are non-negative, we can square both sides and get $(x+y)^2 < 4xy$ so $x^2 + 2xy +y^2 < 4xy$, so $x^2 - 2xy + y^2 < 0$, so $(x-y)^2 < 0$. However, this is a contradiction as no real number squared can be less than 0.
  \end{proof}
\end{enumerate}

\end{document}


%%% Local Variables: 
%%% mode: latex
%%% TeX-master: t
%%% End: 
